\subsection{拓扑模}

定义 3. 给定一个有单位元的拓扑环 A，拓扑左 A-模是一个集合 E，连同：

\begin{enumerate}
\def\labelenumi{\arabic{enumi})}
\item
  一个左 A-模结构；
\item
  一个与 E 的加法群结构相容的拓扑，且满足下列公理：
\end{enumerate}

(MT) 映射 \((\lambda, x) \to \lambda x\)（\(\mathbf{A} \times \mathbf{E}\) 到 \(\mathbf{E}\) 中的映射）是连续的。

我们类似地定义拓扑右 A-模的概念；由于每个右 A-模都可以看作左 A°-模（其中 A° 是 A 的反环），且 A 的拓扑与 A° 的环结构相容，因此无须区分拓扑右 A-模与拓扑左 A°-模。

例. *1) R（相应地，C）上的拓扑向量空间是拓扑 R（相应地，C）模。

\begin{enumerate}
\def\labelenumi{\arabic{enumi})}
\setcounter{enumi}{1}
\tightlist
\item
  设 A 是环，B 是 A 上由 A 的双边理想组成的一个滤子基；设 E 是左 A-模。若给 A 赋以 B 为 o 的基本邻域系的拓扑（与其环结构相容的拓扑）（第 3 节，例 3），并给 E 赋以诸集合 aE（当 a 取遍 B 时）构成 o 的基本邻域系的拓扑（与其加法群结构相容）（§ 1，第 2 节，例），则立即可以验证 E 是拓扑 A-模。
\end{enumerate}

注. 给定拓扑环 A，在左 A-模 E 上考虑一个与 E 的加法群结构相容的拓扑。凭恒等式

\[
\lambda x - \lambda_ {0} x _ {0} = (\lambda - \lambda_ {0}) x _ {0} + \lambda_ {0} (x - x _ {0}) + (\lambda - \lambda_ {0}) (x - x _ {0})
\]

公理 (MT) 等价于下列三条公理的合取：

\((\mathbf{MT}_{\mathbf{I}}^{\prime})\). 对每个 \(x_0\in \mathbf{E}\)，映射 \(\lambda \to \lambda x_0\) 在点 \(\lambda = 0\) 连续

\((\mathbf{MT}_{\mathrm{II}}^{\prime})\). 对每个 \(\lambda_0\in \mathbf{A}\)，映射 \(x\to \lambda_0x\) 在点 \(x = 0\) 连续

\((\mathbf{MT}_{\mathrm{III}}^{\prime})\). 映射 \((\lambda ,x)\to \lambda x\) 在点 (o, o) 连续。

由此我们导出一组充要条件，A-模 E 中 o 的邻域滤子 \(\mathfrak{B}\) 必须满足它们，才能在 E 上定义一个与其模结构相容的拓扑；\(\mathfrak{B}\) 必须满足 § 1 第 2 节的公理 (GA \(_{\mathrm{I}}\) ) 与 (GA \(_{\mathrm{II}}\) )，此外还必须满足下列三条公理：

\((\mathbf{MV}_{\mathrm{I}})\). 对每个 \(x_0 \in E\) 与 V \(\in \mathfrak{B}\)，存在 A 中 o 的一个邻域 S，使得 S. \(x_0 \subset V\)。

\((\mathbf{MV}_{\mathrm{II}})\). 对每个 \(\lambda_0 \in \mathbf{A}\) 与 \(V \in \mathfrak{B}\)，存在 \(W \in \mathfrak{B}\)，使得 \(\lambda_0 W \subset V\)。\((\mathbf{MV}_{\mathrm{III}})\). 对每个 \(V \in \mathfrak{B}\)，存在 \(U \in \mathfrak{B}\) 与一个邻域 \(T\)（\(o\) 在 \(A\) 中的一个邻域），使得 \(T. U \subset V\)。

当环 \(\mathbf{Z}\) 赋以离散拓扑时，每个交换拓扑群都是拓扑 \(\mathbf{Z}\)-模。

若 M 是拓扑 A-模 E 的子模，则显然由 E 的拓扑在 M 上诱导的拓扑与 M 的模结构相容。此外，在商 A-模 E/M 上，E 的拓扑对 M 作出的商拓扑与 A-模结构相容。为看清这一点，只需证明映射 \((\lambda, \dot{x}) \to \lambda \dot{x}\)（\(\mathbf{A} \times (\mathbf{E} / \mathbf{M})\) 到 \(\mathbf{E} / \mathbf{M}\) 上的映射）连续（其中 \(x \to \dot{x}\) 表示 \(\mathbf{E}\) 到 \(\mathbf{E} / \mathbf{M}\) 上的典范映射）。现在，由于我们可以把加法拓扑群 \(\mathbf{A} \times (\mathbf{E} / \mathbf{M})\) 与 \((\mathbf{A} \times \mathbf{E}) / (\{\mathbf{o}\} \times \mathbf{M})\) 等同起来（§ 2，第 9 节，命题 26 的推论），只需证明映射 \((\lambda, x) \to \lambda \dot{x}\)（\(\mathbf{A} \times \mathbf{E}\) 到 \(\mathbf{E} / \mathbf{M}\) 中的映射）连续；而这是直接的，因为所论映射是 \(x \to \dot{x}\) 与 \((\lambda, x) \to \lambda x\) 的复合。

设 \((\mathbf{E}_i)_{i\in \mathbf{I}}\) 是拓扑 A-模的任意一个族，并设 \(\mathbf{E} = \prod_{i\in \mathbf{I}}\mathbf{E}_i\) 是诸 \(\mathbf{E}_i\) 之积这个 A-模。则 E 上的

积拓扑与 E 的 A-模结构相容。为证明这一点，只需证明 A × E 到 E 中的映射 \((\lambda, x) \to (\lambda \cdot \text{pr}_i x)_{i \in I}\) 连续，或者（由第一章 § 4 第 1 节命题 1）对每个指标 \(i \in I\)，映射 \((\lambda, x) \to \lambda \cdot \text{pr}_i x\) 是 A × E 到 \(E_i\) 中的连续映射；但这个映射是 \((\lambda, x_i) \to \lambda x_i\) 与 \((\lambda, x) \to (\lambda, \text{pr}_i x)\) 的复合，而这两个映射都是连续的。

设 A 是豪斯多夫拓扑环，E 是豪斯多夫拓扑 A-模。设 \(\hat{\mathbf{E}}\) 是作为交换拓扑群 E 的完备化的那个加法群（§ 3，第 5 节，定理 2）。Z-双线性映射 \((\lambda, x) \to \lambda x\)（加法群 A、E 之积 \(\mathbf{A} \times \mathbf{E}\) 到加法群 E 中的映射）可由连续性延拓为 \(\hat{\mathbf{A}} \times \hat{\mathbf{E}}\) 到 \(\hat{\mathbf{E}}\) 中的 Z-双线性映射（第 5 节，定理 1），我们仍把这个映射记作 \((\lambda, x) \to \lambda x\)。由恒等延拓原理，我们有 \(\lambda(\mu x) = (\lambda \mu)x\)（对 \(\lambda \in \hat{\mathbf{A}}\)、\(\mu \in \hat{\mathbf{A}}\) 与 \(x \in \hat{\mathbf{E}}\)），且 \(1.x = x\) 对所有 \(x \in \hat{\mathbf{E}}\)；因此外合成法则 \((\lambda, x) \to \lambda x\) 所定义的 \(\hat{\mathbf{A}}\)-模结构在 \(\hat{\mathbf{E}}\) 上与它的拓扑相容。这样定义的拓扑 \(\hat{\mathbf{A}}\)-模 \(\hat{\mathbf{E}}\) 称为拓扑 A-模 E 的完备化。

设 E 是拓扑环 A 上的拓扑模，其中 A 与 E 都不必是豪斯多夫的。设 N（相应地，F）是 \(\{o\}\) 在 A（相应地，E）中的闭包。N 是 A 的双边理想（第 4 节，命题 5），F 是 E 的子 A-模（第 1 节，命题 1）；此外，由连续性我们有 \(\lambda x \in F\)，只要 \(\lambda \in N\) 或 \(x \in F\)。因此我们可以通过过渡到商来定义映射 \((\dot{\lambda}, \dot{x}) \to \dot{\lambda} \dot{x}\)（\((A/N) \times (E/F)\) 到 E/F 中的映射）；容易验证（利用 § 2 第 9 节命题 26 的推论）这个映射连续，从而在 E/F 上定义了一个拓扑 \((A/N)\)-模结构。若令 B = A/N 与 L = E/F，则 B-模 L 称为与 E 相伴的 Hausdorff 模；它的完备化 \(\hat{L}\) 是 \(\hat{A}\)——A 的 Hausdorff 完备化，按定义等于 \(\hat{B}\)——上的一个拓扑模（第 5 节），而这个模 \(\hat{L}\) 称为 E 的 Hausdorff 完备化，记作

\(\hat{\mathbf{E}}\)。与 § 3 第 4 节命题 8 一样可以看到，每个连续同态 \(u: \mathbf{E} \to \mathbf{G}\)（\(\mathbf{E}\) 到完备豪斯多夫 \(\hat{\mathbf{A}}\)-模 \(\mathbf{G}\) 中的连续同态）都唯一分解为 \(u = v \circ \varphi\)，其中 \(v\) 是 \(\hat{\mathbf{E}}\) 到 \(\mathbf{G}\) 中的连续同态，而 \(\varphi\) 是 \(\mathbf{E}\) 到 \(\hat{\mathbf{E}}\) 中的典范映射。我们断言，若 \(\mathbf{E}, \mathbf{E}'\) 是两个拓扑 A-模，且 \(u: \mathbf{E} \to \mathbf{E}'\) 是连续同态，则存在唯一的连续同态 \(\hat{u}: \hat{\mathbf{E}} \to \hat{\mathbf{E}}'\)，使得图表

\[
\begin{array}{c c c} \mathbf {E} & \stackrel {{u}} {{\longrightarrow}} & \mathbf {E} ^ {\prime} \\ \varphi \Big \downarrow & & \Big \downarrow \varphi^ {\prime} \\ \hat {\mathbf {E}} & \stackrel {{\hat {u}}} {{\longrightarrow}} & \hat {\mathbf {E}} ^ {\prime} \end{array}
\]

是交换的，其中 \(\varphi\) 与 \(\varphi'\) 是典范映射。

本节的定义与结果同样适用于任意拓扑环上的伪模，只需删去所有关于环的单位元的提述。

